\begin{lemma}
Let $n\ge1$ be an integer, let $1<p<\infty$, and let
$k\in\mathbb{Z}$.  Let $a\colon\mathbb{R}\to\mathbb{C}$ be continuous on
$[0,1)$, $C^\infty$ on $(0,1)$, and a solution of the Euler equation
\[
  r\bigl(r\,a'\bigr)'(r)+pn\,r\,a'(r)+(n^2-k^2)(p-1)\,a(r)=0
  \qquad\text{for every }r\in(0,1).
\]
Then
\[
  a(r)=\begin{cases}
    0 & \text{if }|k|<n,\\
    A\,r^{\lambda_{|k|}(n,p)} & \text{if }|k|\ge n,
  \end{cases}
  \qquad\text{for every }r\in[0,1),
\]
where in the second case $A\in\mathbb{C}$ is a constant depending on $a$, and
$\lambda_{|k|}(n,p)$ is the exponent of \noderef{LambdaCoef}.  Real powers are
taken with the conventions $0^{\lambda}=0$ for $\lambda>0$ and $0^{0}=1$.

This is formula (36) of the paper, with $\lambda_k=\lambda_{|k|}(n,p)$ its (37);
the two agree because only $k^2$ enters the formula, which is the paper's remark
$\lambda_k=\lambda_{-k}$ at line 331.  No restriction $p\neq2$ is imposed: the
conclusion holds for every $1<p<\infty$, including the single pair $(k,p)=(0,2)$
at which the indicial roots coincide.
\end{lemma}
\begin{proof}
Write $S=\sqrt{4k^2(p-1)+n^2(p-2)^2}$ and
\[
  \lambda^{+}=\lambda_{|k|}(n,p)=\frac{-pn+S}{2},\qquad
  \lambda^{-}=-pn-\lambda^{+}=\frac{-pn-S}{2},
\]
as in \noderef{LambdaCoef} (using $|k|^2=k^2$).

\emph{Step 0a: the two cases.}  Since $p>1$ we have $4k^2(p-1)\ge0$, with
strict inequality when $k\neq0$; and $n^2(p-2)^2\ge0$, with strict inequality
when $p\neq2$, because $n\ge1$.  Hence the radicand
$4k^2(p-1)+n^2(p-2)^2$ is positive unless $k=0$ and $p=2$, in which case it is
$0$.  We therefore distinguish the \emph{non-degenerate case} ($k\neq0$ or
$p\neq2$), treated in Steps 0--5, and the \emph{degenerate case} ($k=0$ and
$p=2$), treated in Step 6.  The degenerate case is the one in which the paper's
assertion at line 313 that the roots are distinct fails.  Only the bound (4) of
Step 0b is used in both cases.

\emph{Step 0b: $a$ is bounded near $0$; this is common to both cases.}  The
interval $[0,\tfrac12]$ is a compact subset of $[0,1)$ and $a$ is continuous on
$[0,1)$, so
\begin{equation}\tag{4}
  M=\sup_{r\in[0,1/2]}|a(r)|<\infty .
\end{equation}

\emph{Steps 0--5 assume $k\neq0$ or $p\neq2$.}  Then the radicand is positive
by Step 0a, so
\begin{equation}\tag{1}
  S>0,\qquad S^2=4k^2(p-1)+n^2(p-2)^2,\qquad \lambda^{+}-\lambda^{-}=S>0 .
\end{equation}

\emph{Step 0: $\lambda^{-}<0$ always, and the sign of $\lambda^{+}$ in three
regimes.}  Since $n\ge1$ and $p>1$ we have $pn>0$, and $S\ge0$, so
\begin{equation}\tag{2}
  \lambda^{-}=\frac{-pn-S}{2}<0 .
\end{equation}
For $\lambda^{+}$, note that $\lambda^{+}$ and $0$ compare exactly as $S$ and
$pn$ do, by (1); and both $S\ge0$ and $pn>0$, so $S$ and $pn$ compare exactly as
$S^2$ and $(pn)^2$ do.  Now, using the identity $p^2-(p-2)^2=4(p-1)$,
\[
  S^2-(pn)^2=4k^2(p-1)+n^2(p-2)^2-n^2p^2
  =4k^2(p-1)-n^2\bigl(p^2-(p-2)^2\bigr)=4(p-1)\bigl(k^2-n^2\bigr),
\]
and $4(p-1)>0$.  Therefore
\begin{equation}\tag{3}
  \lambda^{+}<0\iff|k|<n,\qquad \lambda^{+}=0\iff|k|=n,\qquad \lambda^{+}>0\iff|k|>n,
\end{equation}
since $k^2-n^2$ is negative, zero or positive exactly as $|k|$ is less than,
equal to or greater than $n$.  The middle regime is the one the paper's
sentence at line 323 passes over: at $|k|=n$ the exponent is exactly $0$, so
$r^{\lambda^{+}}$ is the constant $1$, which is bounded, and the corresponding
constant need \emph{not} vanish.  Note also that the roots remain distinct
there, since $\lambda^{-}<0=\lambda^{+}$ by (2).

\emph{Step 2: the two constants.}  The hypotheses of
\noderef{EulerOdeSolution} are met: $n\ge1$, $1<p<\infty$, its non-degeneracy
hypothesis ``$k\neq0$ or $p\neq2$'' is the case assumption in force, $a$ is
$C^\infty$ on $(0,1)$, and $a$ satisfies the Euler equation there.  It
therefore provides $A^{+},A^{-}\in\mathbb{C}$ with
\begin{equation}\tag{5}
  a(r)=A^{+}r^{\lambda^{+}}+A^{-}r^{\lambda^{-}}\qquad\text{for every }r\in(0,1).
\end{equation}

\emph{Step 3: $A^{-}=0$.}  Multiply (5) by $r^{-\lambda^{-}}$, which is
legitimate for $r>0$, and use $r^{-\lambda^{-}}r^{\lambda^{+}}=r^{\lambda^{+}-\lambda^{-}}=r^{S}$
and $r^{-\lambda^{-}}r^{\lambda^{-}}=1$:
\begin{equation}\tag{6}
  r^{-\lambda^{-}}a(r)=A^{+}r^{S}+A^{-}\qquad\text{for every }r\in(0,1).
\end{equation}
Let $r\to0^{+}$.  By (2), $-\lambda^{-}>0$, so $r^{-\lambda^{-}}\to0$, and by
(4), $|r^{-\lambda^{-}}a(r)|\le M\,r^{-\lambda^{-}}\to0$ for $r\in(0,\tfrac12]$;
thus the left-hand side of (6) tends to $0$.  By (1), $S>0$, so
$r^{S}\to0$ and the first term on the right tends to $0$.  Hence $A^{-}=0$, and
(5) becomes
\begin{equation}\tag{7}
  a(r)=A^{+}r^{\lambda^{+}}\qquad\text{for every }r\in(0,1).
\end{equation}

\emph{Step 4: the case $|k|<n$.}  Here $\lambda^{+}<0$ by (3).  Suppose
$A^{+}\neq0$.  Then by (7), for $r\in(0,\tfrac12]$,
$|a(r)|=|A^{+}|r^{\lambda^{+}}\to\infty$ as $r\to0^{+}$, because
$-\lambda^{+}>0$; this contradicts the bound (4).  Hence $A^{+}=0$, and by (7)
$a(r)=0$ for every $r\in(0,1)$.  Finally $a$ is continuous at $0$ from the
right, so $a(0)=\lim_{r\to0^{+}}a(r)=0$ as well.  Therefore $a(r)=0$ for every
$r\in[0,1)$, which is the first case of the assertion.

\emph{Step 5: the case $|k|\ge n$.}  Here $\lambda^{+}\ge0$ by (3).  Put
$A=A^{+}$.  By (7), $a(r)=A\,r^{\lambda^{+}}$ for every $r\in(0,1)$, so it
remains only to check the equality at $r=0$, i.e.\ that
$a(0)=A\cdot0^{\lambda^{+}}$.  We distinguish the two sub-regimes of (3).

If $|k|>n$ then $\lambda^{+}>0$, so $0^{\lambda^{+}}=0$ by the convention for
real powers, and $r^{\lambda^{+}}\to0$ as $r\to0^{+}$; by continuity of $a$ at
$0$ and (7),
$a(0)=\lim_{r\to0^{+}}A\,r^{\lambda^{+}}=0=A\cdot0^{\lambda^{+}}$.

If $|k|=n$ then $\lambda^{+}=0$, so $r^{\lambda^{+}}=1$ for every $r>0$ and
$0^{\lambda^{+}}=0^{0}=1$.  Thus (7) says $a(r)=A$ for every $r\in(0,1)$, and by
continuity of $a$ at $0$, $a(0)=\lim_{r\to0^{+}}A=A=A\cdot0^{\lambda^{+}}$.  This
is the regime in which the coefficient is a nonzero constant; it is exactly the
$k=\pm n$ branch of (36), and it is consistent with the paper's remark
$\lambda_n=0$ at line 331.

In both sub-regimes $a(r)=A\,r^{\lambda^{+}}$ for every $r\in[0,1)$, which is
the second case of the assertion.  This completes the non-degenerate case.

\emph{Step 6: the degenerate case $k=0$, $p=2$.}  Here only the first branch of
the assertion has to be proved, and the second is vacuous: $|k|=0$, and
$0<n$ because $n\ge1$, so $|k|<n$ and the hypothesis $|k|\ge n$ of the second
branch is never met.

With $k=0$ and $p=2$ the Euler equation in the hypothesis reads
\[
  r\bigl(r\,a'\bigr)'(r)+2n\,r\,a'(r)+n^2\,a(r)=0\qquad\text{for every }r\in(0,1),
\]
because $pn=2n$ and $(n^2-k^2)(p-1)=n^2\cdot1=n^2$.  Together with $n\ge1$ and
the $C^\infty$-smoothness of $a$ on $(0,1)$, these are exactly the hypotheses of
\noderef{EulerOdeDegenerate}, which therefore provides $A,B\in\mathbb{C}$ with
\begin{equation}\tag{8}
  a(r)=\bigl(A+B\log r\bigr)r^{-n}\qquad\text{for every }r\in(0,1).
\end{equation}

We show $B=0$.  Suppose $B\neq0$.  For $0<r\le\tfrac12$ we have $r^{-n}\ge1$,
since $n\ge1$ and $0<r\le1$; and $\log r<0$, so $|\log r|=-\log r$.  Hence by
(8) and the triangle inequality, for $0<r\le\tfrac12$,
\[
  |a(r)|=\bigl|A+B\log r\bigr|\,r^{-n}
  \ge\bigl(|B|\,|\log r|-|A|\bigr)r^{-n}\ge|B|\,|\log r|-|A|
\]
whenever the middle quantity is nonnegative.  As $r\to0^{+}$ we have
$|\log r|\to\infty$, so $|B||\log r|-|A|\to\infty$, contradicting the bound
$|a(r)|\le M$ of (4).  Therefore $B=0$.

Now (8) reads $a(r)=A\,r^{-n}$ on $(0,1)$.  If $A\neq0$ then
$|a(r)|=|A|r^{-n}\to\infty$ as $r\to0^{+}$, because $n\ge1$, again
contradicting (4).  Therefore $A=0$ and $a(r)=0$ for every $r\in(0,1)$.
Finally $a$ is continuous at $0$ from the right, so
$a(0)=\lim_{r\to0^{+}}a(r)=0$ as well, and $a(r)=0$ for every $r\in[0,1)$.
This is the first case of the assertion, and it is the only one that can arise
in the degenerate case.
\end{proof}
